\section{Mode Contractions in the Abstract State}
\label{sec:ch19-mode-correlators}

The preceding section constructed a normalized positive functional on an
\emph{uncompressed} oscillator algebra. We now ask a more concrete question:
what does that state predict for the product of two density modes? A
Bogoliubov transformation mixes an oscillator's creation operator on one
branch with an annihilation operator on the opposite branch. Consequently,
even if the transformed (dressed) modes define a vacuum, the original
(\emph{bare}) modes need not annihilate it. The outcome is a nonzero
correlation between pairs of bare operators. This is the elementary
algebraic mechanism behind the mode contractions derived below.

There is a crucial choice before computing anything: does the observable
contain $AB$ or $BA$? These are different measurements whenever $A$ and
$B$ do not commute, and a state is a \emph{linear}, not a multiplicative,
evaluator. In particular, exchanging the order of a bosonic creation and
annihilation mode changes its vacuum expectation by their scalar
commutator. We therefore compute \emph{ordered} contractions first; the
symmetrized combinations will be identified only afterward. This section
works solely in the abstract oscillator functional of
Section~\ref{sec:ch19-scope}. It does not evaluate a ground-state vector
or density matrix in the finite fermionic energy-budget space.

\subsection{Mode labels, hypotheses, and the vacuum rule}

Fix a nonnegative integer $M$. The set of positive mode labels is
\[
  \mathcal{I}_M=\{1,2,\ldots,M\},
  \qquad \nu\in\{R,L\}.
\]
When $M=0$ the statements about individual modes are empty; whenever we
choose a particular mode below we assume $M\geq1$. The formalized index
$i\in\operatorname{Fin}(M)$ corresponds to $m=i+1$, so every oscillator
weight $m$ is strictly positive. The lattice length $L$ is reserved for
the spatial Fourier sums of the next section (in the supplied finite-lattice
implementation $L=2h$); it is \emph{not} the dimension of the oscillator
polynomial carrier. The mode algebra considered here has $2M$ oscillator
species but admits arbitrarily high occupation in each species.

Write $C_{\nu m}$ for the bare creation generator and
$A_{\nu m}:=C_{\nu m}^{*}$ for its adjoint, the annihilation generator.
We use the density-mode correspondence of the source convention,
\[
  \rho_{m,\nu}=C_{\nu m},\qquad
  \rho_{-m,\nu}=A_{\nu m}\quad(m\in\mathcal I_M).
\]
The exact uncompressed canonical commutation relations (CCR) are
\begin{align}
 [A_{\nu m},C_{\mu n}]&=m\,\delta_{\nu\mu}\delta_{mn}\,I,
   \label{eq:ch19-modes-ccr}\\
 [A_{\nu m},A_{\mu n}]&=0,\qquad
 [C_{\nu m},C_{\mu n}]=0.
\end{align}
Here $\delta_{mn}$ is one if $m=n$ and zero otherwise. Operators on
different branches commute, but, as we shall see, they can still have
nonzero expectations of their products in a correlated state.

Let $\omega_0$ denote the polynomial vacuum functional defined in
Section~\ref{sec:ch19-scope}. In its explicit polynomial representation,
$C_{\nu m}$ multiplies by the variable $X_{\nu m}$, while
$A_{\nu m}=m\,\partial/\partial X_{\nu m}$. The vacuum is the
constant polynomial $1$. Therefore
\begin{align}
 A_{\nu m}1&=0,
 &C_{\nu m}1&=X_{\nu m},
 &A_{\nu m}X_{\mu n}
   &=m\delta_{\nu\mu}\delta_{mn}\,1.
\end{align}
Applying the operators from right to left, then extracting the vacuum
coefficient, yields
\begin{align}
 \omega_0(A_{\nu m}C_{\mu n})
   &=m\delta_{\nu\mu}\delta_{mn},
   &\omega_0(C_{\nu m}A_{\mu n})&=0,
     \label{eq:ch19-free-two-orders}\\
 \omega_0(A_{\nu m}A_{\mu n})&=0,
   &\omega_0(C_{\nu m}C_{\mu n})&=0.
     \label{eq:ch19-free-like-pairs}
\end{align}
For example, $A_{\nu m}C_{\nu m}1=m1$, whereas
$C_{\nu m}A_{\nu m}1=0$. This two-line calculation is the source of
the distinction between the $c^2$ and $s^2$ contractions. The identities
are exact on the uncompressed polynomial carrier: no top-occupation
boundary has been discarded.

Let $c,s\in\mathbb R$ satisfy
\begin{equation}
  c^2-s^2=1.
  \label{eq:ch19-hyperbolic-condition}
\end{equation}
This is the \emph{real hyperbolic-pair} hypothesis of Chapter~18; it
includes the choice $c=1,s=0$. No sign condition on $s$ is necessary.
The Bogoliubov generator map $\beta=\beta_{c,s}$ acts by
\begin{align}
 \beta(C_{Rm})&=cC_{Rm}+sA_{Lm},
 &\beta(A_{Rm})&=cA_{Rm}+sC_{Lm},\label{eq:ch19-beta-right}\\
 \beta(C_{Lm})&=cC_{Lm}+sA_{Rm},
 &\beta(A_{Lm})&=cA_{Lm}+sC_{Rm}.
 \label{eq:ch19-beta-left}
\end{align}
Because the coefficients are real, the second column of each line is
the adjoint of the first. Because $c^2-s^2=1$, applying the same map
with $s$ replaced by $-s$ gives its inverse. For example,
\[
 \beta^{-1}(C_{Rm})=cC_{Rm}-sA_{Lm},\qquad
 \beta^{-1}(A_{Rm})=cA_{Rm}-sC_{Lm}.
\]
The other branch follows by interchanging $R$ and $L$. These identities
are a special case of the invertible generator transformation established
in Chapter~18 and in the supplied A08 development.

The \emph{dressed} modes are $\widetilde C_{\nu m}=\beta(C_{\nu m})$
and $\widetilde A_{\nu m}=\beta(A_{\nu m})$. As established in the
preceding section, the dressed state is the \emph{inverse pullback}
\begin{equation}
  \omega_P(B)=\omega_0\bigl(\beta^{-1}(B)\bigr),
  \qquad B\in\mathfrak A_{\mathrm{CCR}}.
  \label{eq:ch19-pullback}
\end{equation}
This direction is essential. In fact,
$\omega_P(B\widetilde A_{\nu m})=0$ and
$\omega_P(\widetilde C_{\nu m}B)=0$ for every polynomial observable
$B$, since applying $\beta^{-1}$ reduces these expressions to the
undressed vacuum rules. The shorthand $\omega_{\widetilde\Omega}$ in
the original chapter notes refers to this same abstract functional;
we use $\omega_P$ consistently here. In particular, the equations below
are not assertions about finite-CAR expectations
$\langle\Psi,B\Psi\rangle$ or $\operatorname{Tr}(\varrho B)$.

\subsection{Same-branch contractions: why both orders survive}

\begin{theorem}[Both ordered right- and left-branch contractions]
\label{thm:ch19-mode-same-branch}
For any $M\in\mathbb N$, real $c,s$ satisfying
\eqref{eq:ch19-hyperbolic-condition}, $m,n\in\mathcal I_M$, and
$\nu\in\{R,L\}$, the functional~\eqref{eq:ch19-pullback} obeys
\begin{align}
 \omega_P(C_{\nu m}A_{\nu n})
   &=m\delta_{mn}s^2,
   \label{eq:ch19-mode-CA}\\
 \omega_P(A_{\nu m}C_{\nu n})
   &=m\delta_{mn}c^2,
   \label{eq:ch19-mode-AC}\\
 \omega_P(C_{\nu m}C_{\nu n})
   &=0,\qquad
 \omega_P(A_{\nu m}A_{\nu n})=0.
 \label{eq:ch19-mode-like}
\end{align}
All products are written in their actual evaluation order.
\end{theorem}

\begin{proof}
It suffices to treat the right branch; the opposite-branch exchange
$R\leftrightarrow L$ preserves the CCR and the transformation.
Multiplicativity of the \emph{algebra map} $\beta^{-1}$ (not of the
state $\omega_0$) gives
\begin{align*}
 \omega_P(C_{Rm}A_{Rn})
 &=\omega_0\bigl[(cC_{Rm}-sA_{Lm})
                     (cA_{Rn}-sC_{Ln})\bigr]\\
 &=c^2\omega_0(C_{Rm}A_{Rn})
   -cs\omega_0(C_{Rm}C_{Ln})\\
 &\quad-cs\omega_0(A_{Lm}A_{Rn})
   +s^2\omega_0(A_{Lm}C_{Ln})\\
 &=s^2m\delta_{mn}.
\end{align*}
Every vanishing term has been removed using
\eqref{eq:ch19-free-two-orders}--\eqref{eq:ch19-free-like-pairs}.
For the reversed product the expansion is different:
\begin{align*}
 \omega_P(A_{Rm}C_{Rn})
 &=\omega_0\bigl[(cA_{Rm}-sC_{Lm})
                     (cC_{Rn}-sA_{Ln})\bigr]\\
 &=c^2\omega_0(A_{Rm}C_{Rn})
   -cs\omega_0(A_{Rm}A_{Ln})\\
 &\quad-cs\omega_0(C_{Lm}C_{Rn})
   +s^2\omega_0(C_{Lm}A_{Ln})\\
 &=c^2m\delta_{mn}.
\end{align*}
The $CC$ product is likewise mapped to
\begin{align*}
 \omega_0[(cC_{Rm}-sA_{Lm})(cC_{Rn}-sA_{Ln})]
 &=c^2\omega_0(C_{Rm}C_{Rn})\\
 &\quad-cs\omega_0(C_{Rm}A_{Ln})
   -cs\omega_0(A_{Lm}C_{Rn})\\
 &\quad+s^2\omega_0(A_{Lm}A_{Ln})=0,
\end{align*}
since in the potentially surviving $AC$ term the branches differ.
The $AA$ identity follows by adjunction or by the same four-term
expansion. Nothing in the proof requires a finite-dimensional
representation of the CCR.
\end{proof}

Subtracting the first contraction from the second gives the exact
commutator check
\begin{equation}
 \omega_P(A_{\nu m}C_{\nu n})
 -\omega_P(C_{\nu n}A_{\nu m})
 =m\delta_{mn}(c^2-s^2)=m\delta_{mn}.
 \label{eq:ch19-mode-commutator-check}
\end{equation}
The right-hand side is fixed by the algebra, independently of the
chosen state. Thus the coefficient $s^2$ cannot simply replace
$c^2$ in the reversed contraction. More physically, even in a free
vacuum an annihilator followed by a creator can produce a nonzero
fluctuation: the creator first excites the vacuum, then the annihilator
returns it to the vacuum.

\subsection{Cross-branch contractions and the sign of the mixing}

The transformation also couples opposite branches. Correlation does not
require a nonzero cross-branch commutator: here the $R$ and $L$
oscillators commute, but $\omega_P$ is not a product state of the
\emph{bare} right- and left-branch vacuum functionals. The distinction
between algebraic commutativity and statistical independence is
important in interpreting a Bogoliubov-transformed ground state.

\begin{theorem}[Opposite-branch anomalous contractions]
\label{thm:ch19-mode-cross-branch}
Under the hypotheses of Theorem~\ref{thm:ch19-mode-same-branch},
\begin{align}
 \omega_P(C_{Rm}C_{Ln})
   &=-m\delta_{mn}cs,
   \label{eq:ch19-cross-CC}\\
 \omega_P(A_{Rm}A_{Ln})
   &=-m\delta_{mn}cs,
   \label{eq:ch19-cross-AA}\\
 \omega_P(C_{Rm}A_{Ln})
   &=0,\qquad
 \omega_P(A_{Rm}C_{Ln})=0.
   \label{eq:ch19-cross-mixed}
\end{align}
The same results hold after exchanging $R$ and $L$.
\end{theorem}

\begin{proof}
Consider first the product of two creators. Applying the inverse
transformation and expanding without commuting factors gives
\begin{align*}
 \omega_P(C_{Rm}C_{Ln})
 &=\omega_0\bigl[(cC_{Rm}-sA_{Lm})
                     (cC_{Ln}-sA_{Rn})\bigr]\\
 &=c^2\omega_0(C_{Rm}C_{Ln})
   -cs\omega_0(C_{Rm}A_{Rn})\\
 &\quad-cs\omega_0(A_{Lm}C_{Ln})
   +s^2\omega_0(A_{Lm}A_{Rn})\\
 &=-cs\,m\delta_{mn}.
\end{align*}
The single surviving term has precisely the minus sign introduced
by $\beta^{-1}$; using the forward map instead would reverse this
sign and would describe a different state.
For two annihilators the analogous computation is
\begin{align*}
 \omega_P(A_{Rm}A_{Ln})
 &=\omega_0\bigl[(cA_{Rm}-sC_{Lm})
                     (cA_{Ln}-sC_{Rn})\bigr]\\
 &=c^2\omega_0(A_{Rm}A_{Ln})
   -cs\omega_0(A_{Rm}C_{Rn})\\
 &\quad-cs\omega_0(C_{Lm}A_{Ln})
   +s^2\omega_0(C_{Lm}C_{Rn})\\
 &=-cs\,m\delta_{mn}.
\end{align*}
Finally, the mixed-branch product becomes
\begin{align*}
 \omega_P(C_{Rm}A_{Ln})
 &=c^2\omega_0(C_{Rm}A_{Ln})
   -cs\omega_0(C_{Rm}C_{Rn})\\
 &\quad-cs\omega_0(A_{Lm}A_{Ln})
   +s^2\omega_0(A_{Lm}C_{Rn})=0.
\end{align*}
The other mixed-branch contraction vanishes by the same vacuum
identities. Exchanges of branches are permitted because they commute
in the CCR algebra. This completes the proof.
\end{proof}

These cross-branch quantities are often called \emph{anomalous}
contractions: they pair two creators or two annihilators, which would
have zero expectation in the undressed polynomial vacuum. Their
nonzero value in $\omega_P$ encodes the pair mixing; the sign records
our choice of the inverse-pullback convention and the sign of $s$.

\subsection{Ordered versus symmetrized expectations}

For any two algebra elements $B,D$, introduce their symmetrized
product $\{B,D\}:=BD+DB$. The ordered and symmetrized functionals
measure different algebra elements:
\[
 \omega_P(BD)
 =\frac12\omega_P(\{B,D\})
  +\frac12\omega_P([B,D]).
\]
For one right- or left-branch creator/annihilator pair, the previous
theorem implies
\begin{equation}
 \frac12\omega_P(\{A_{\nu m},C_{\nu n}\})
   =\frac{m\delta_{mn}}2(c^2+s^2).
 \label{eq:ch19-mode-anticommutator}
\end{equation}
The antisymmetric half is
$\frac12\omega_P([A_{\nu m},C_{\nu n}])=
\frac12m\delta_{mn}$. Thus symmetrization keeps the sum of the
fluctuations but suppresses the information about which ordered
product carried the vacuum term. The anomalous $CC$ and $AA$
contractions do not have this ordering ambiguity across branches:
the corresponding generators commute.

It is useful to express these coefficients in the parameter convention
introduced in Chapter~18. Put
\begin{equation}
 g=(c-s)^2>0,\qquad g^{-1}=(c+s)^2.
 \label{eq:ch19-mode-g}
\end{equation}
The inverse equality follows from $(c-s)(c+s)=1$. Hence
\begin{align}
 c^2+s^2&=\frac{g+g^{-1}}2,
 & c^2&=\frac{g+g^{-1}+2}{4},
 & s^2&=\frac{g+g^{-1}-2}{4},\\
 cs&=\frac{g^{-1}-g}{4}.
 \label{eq:ch19-mode-g-coefficients}
\end{align}
In particular,
\begin{align*}
 \frac12\omega_P(\{A_{\nu m},C_{\nu n}\})
   &=\frac{m\delta_{mn}}4(g+g^{-1}),\\
 \omega_P(C_{Rm}C_{Ln})
   &=-\frac{m\delta_{mn}}4(g^{-1}-g).
\end{align*}
The coupling $g$ here is a \emph{hyperbolic/Luttinger parameter}; it
is not the Umklapp coupling $g_U$ used in a later chapter.
These finite algebraic identities do not invoke a continuum limit.

\subsection{Worked example: two modes and a squeezed combination}

Choose $M=2$ and the real hyperbolic pair
\[
 c=\frac53,\qquad s=\frac43,\qquad
 c^2-s^2=\frac{25-16}{9}=1,
 \qquad g=(c-s)^2=\frac19.
\]
At $m=n=1$ the right-branch ordered contractions are
\[
 \omega_P(C_{R1}A_{R1})=\frac{16}{9},\qquad
 \omega_P(A_{R1}C_{R1})=\frac{25}{9}.
\]
Their difference is exactly $1$, as required by the uncompressed
commutator. The cross-branch anomalous contractions are
\[
 \omega_P(C_{R1}C_{L1})
 =\omega_P(A_{R1}A_{L1})=-\frac{20}{9}.
\]
For comparison, mode orthogonality yields
$\omega_P(C_{R1}A_{R2})=\omega_P(C_{R1}C_{L2})=0$.
At the second mode the nonzero contractions simply acquire the
weight $m=2$:
\[
 \omega_P(C_{R2}A_{R2})=\frac{32}{9},\quad
 \omega_P(A_{R2}C_{R2})=\frac{50}{9},\quad
 \omega_P(C_{R2}C_{L2})=-\frac{40}{9}.
\]

There is also an instructive positive quadratic observable. Fix one
mode $m$ and define $Y_m=C_{Rm}+A_{Lm}$. Its adjoint is
$Y_m^*=A_{Rm}+C_{Lm}$. Expanding $Y_m^*Y_m$ in precisely this order,
\begin{align*}
 \omega_P(Y_m^*Y_m)
 &=\omega_P(A_{Rm}C_{Rm})
   +\omega_P(A_{Rm}A_{Lm})\\
 &\quad+\omega_P(C_{Lm}C_{Rm})
   +\omega_P(C_{Lm}A_{Lm})\\
 &=m(c^2-2cs+s^2)=m(c-s)^2=mg.
\end{align*}
In our numerical example, $\omega_P(Y_1^*Y_1)=1/9>0$.
Although individual contributions include negative cross-branch
correlations, the positive functional still gives a nonnegative
expectation to an element of the form $Y^*Y$. This calculation makes
the underlying tradeoff visible: the state can encode correlations
that are simple in the full oscillator representation, while a
finite energy-budget compression cannot be assumed to preserve these
CCR-based manipulations without additional boundary analysis.

\subsection{Free-pair limit and exact scope of the result}

The supplied implementation chooses the free hyperbolic pair
$(c,s)=(1,0)$, so $\beta$ and $\beta^{-1}$ are the identity and
$\omega_P=\omega_0$. The complete nonzero two-point data then reduce
to
\begin{equation}
 \omega_0(A_{\nu m}C_{\mu n})
 =m\delta_{\nu\mu}\delta_{mn},
 \label{eq:ch19-mode-free-limit}
\end{equation}
while $\omega_0(C_{\nu m}A_{\mu n})=0$ and all $CC$, $AA$
and opposite-branch contractions vanish. Thus the normal-ordered
free-vacuum occupation is zero, but the reverse-ordered vacuum
fluctuation is \emph{not} zero. Omitting it would make the subsequent
spatial density correlator incorrect already in the free theory.

All results in this section are \emph{exact algebraic evaluations} for
finitely many positive mode labels in the positive abstract state
$\omega_P$. They follow from the polynomial vacuum rules, the real
hyperbolic condition, and the inverse Bogoliubov pullback. In the
supplied Lean sources these evaluations correspond to the same-branch
and cross-branch contraction theorems in
	exttt{BosonizeStubs/Ch19Correlations.lean}, together with its separate
vanishing lemmas. Their acceptance in the Chapter~19/A08
Phase~B development does not constitute promotion to the frozen Core.

The next section will combine the ordered contractions with a
\emph{specified} finite Fourier orientation and normalization to
construct spatial density kernels. That step must retain both the
$c^2$ and $s^2$ terms; replacing the ordered product by a cosine-only
symmetrization changes the observable. A further, independent task is
to compare these abstract values with expectations in a constructed
finite fermionic ground state, including any projection or boundary
residuals. No such equality is assumed here.
